% =============================================================================
% Project : THE_PHILOSOPHY_OF_VIABILITY
% Book    : Book I — Ontology of Viability
% Book ID : VS.B01
%
% File:
% VS.B01_ONTOLOGY_OF_VIABILITY/2_MAINMATTER/
% PART_I_IDENTITY_AND_MANIFESTATION/
% CHAPTER_07_FORMAL_AND_LOGICAL_RECONSTRUCTION/
% sections/section_06_formalization_of_the_principles_of_identity.tex
%
% Section:
% Chapter 07, Section 06 — Formalization of the Principles of Identity
%
% Structural Location:
% VS.B01.P01.C07.S06
%
% Purpose:
% To develop provisional formal schemas for the three foundational
% commitments concerning Being and Identity, while preserving their
% distinct logical contents and avoiding unintended ontological
% commitments introduced by formal representation.
%
% =============================================================================

\section{Formalization of the Principles of Identity}
\label{sec:ch07_s06}

\subsection{The Task of Formalizing Foundational Principles}

The preceding section identified three foundational commitments
governing the relation between Being and Identity:

\begin{enumerate}
    \item Analytical Distinction;
    \item Ontological Inseparability;
    \item Absence of Ontological Priority.
\end{enumerate}

The present task is to examine how these commitments may be
represented within a formal framework.

Formalization requires more than translating philosophical
sentences into symbolic expressions. A formal representation
must identify the logical structure of a claim, specify the
domain over which it applies, and make explicit the relations
that its symbols are intended to represent.

The central difficulty is that the three commitments concern
different dimensions of the relation between Being and Identity.
A formalization that represents them through a single equivalence,
a single relation, or an undifferentiated conjunction would risk
concealing their independent contents.

The formal reconstruction must therefore preserve three
distinct questions:

\begin{itemize}
    \item Can Being and Identity be distinguished analytically?
    \item Are Being and Identity ontologically inseparable?
    \item Does either possess ontological priority over the other?
\end{itemize}

These questions require different formal resources. Their
representations must be connected without being collapsed
into one another.

The formal schemas developed below are provisional. They
identify the logical requirements of each commitment and
provide a basis for subsequent examination of their consistency,
independence, and consequences. They do not claim that a
complete formal language for the ontology has already been
established.

\subsection{Preliminary Formal Vocabulary}

A formal representation requires a minimal vocabulary
capable of distinguishing an existent from the analytical
and ontological dimensions through which it is examined.

Let \(D_E\) denote a domain of discourse whose members are
the existents under consideration.

Let the following expressions be used provisionally:

\begin{itemize}
    \item \(E(x)\): \(x\) is an existent within the domain
    under consideration;

    \item \(B(x)\): the dimension of Being is applicable
    to the actuality of \(x\);

    \item \(I(x)\): the dimension of Identity is applicable
    to \(x\);

    \item \(A(B,I)\): Being and Identity are analytically
    distinguishable;

    \item \(S_O(B,I)\): Being and Identity are ontologically
    inseparable;

    \item \(P_O(B,I)\): Being has ontological priority
    over Identity;

    \item \(P_O(I,B)\): Identity has ontological priority
    over Being.
\end{itemize}

The expressions \(B(x)\) and \(I(x)\) must be interpreted
with care. They are formal indicators of the applicability
of the respective dimensions in the analysis of an existent.
They are not definitions of Being and Identity as two
independent entities, nor do they establish that Identity
is an additional property attached to an already constituted
existent.

Likewise, \(A\), \(S_O\), and \(P_O\) are provisional
meta-level predicates used to represent relations between
theoretical dimensions. Their introduction does not
presuppose that Being and Identity are members of a
common ontological category.

The formal vocabulary is therefore an instrument of
representation. Its interpretation must remain governed
by the ontological commitments it is intended to express.

\subsection{Formalization of Analytical Distinction}

\subsubsection{The Content of Analytical Distinction}

Analytical Distinction states that Being and Identity
can be examined through different conceptual questions
and formulations.

The distinction makes it possible to ask what it means
for an existent to be actual and what it means for that
existent to be the existent that it is.

These questions are related, but they are not identical.
The first concerns the actuality of what exists; the
second concerns the identity of the existent.

The formal task is to represent this difference without
turning it into a claim that Being and Identity are
ontologically independent.

\subsubsection{A Provisional Formal Schema}

Let \(\mathcal{Q}\) denote a domain of admissible analytical
questions or conceptual roles.

Let \(q_B\) designate a question directed toward Being,
and \(q_I\) a question directed toward Identity.

Analytical distinction may then be represented by the
following schema:

\[
\mathsf{AD}:
\quad
q_B \neq q_I
\quad \land \quad
\mathsf{Content}(q_B)
\neq
\mathsf{Content}(q_I)
\]

The notation expresses that the two analytical questions
are not interchangeable and that they have different
conceptual contents.

The schema is not a claim that two independent entities
called Being and Identity must exist. It represents
a distinction between analytical roles.

A more explicit formulation may be stated as:

\[
\mathsf{AD}:
\quad
\exists q_B,q_I \in \mathcal{Q}
\left[
\mathsf{DirectedTo}(q_B,B)
\land
\mathsf{DirectedTo}(q_I,I)
\land
q_B \neq q_I
\right]
\]

This formulation identifies the difference between
the two analytical directions. It does not yet
provide a complete semantic theory of the terms
Being and Identity.

\subsubsection{What the Formalization Does Not Establish}

The distinction between \(q_B\) and \(q_I\) does not
establish that Being and Identity can exist separately.

In particular, the following inference is invalid:

\[
q_B \neq q_I
\quad \therefore \quad
B \text{ and } I
\text{ are ontologically separable}
\]

The conclusion does not follow from the premise.
A difference between analytical questions does not
establish a corresponding separation in reality.

Nor does the schema establish that Being and Identity
are two properties of an existent. Their formal
representation through different predicates or
analytical indicators must not be mistaken for
a commitment to a property-based ontology.

The distinction is therefore analytical in its
intended interpretation. Its ontological consequences
must be assessed through the independent commitment
of Ontological Inseparability.

\subsection{Formalization of Ontological Inseparability}

\subsubsection{The Content of Ontological Inseparability}

Ontological Inseparability states that Being and
Identity are inseparable in the actuality of an
existent.

An existent is not first constituted as an actuality
that lacks identity and then supplemented by
identity as a subsequent addition.

Likewise, Identity is not treated as an independently
constituted ontological item that can stand apart
from the existent whose identity it is.

The formalization must express the joint applicability
of Being and Identity to an existent while preserving
the fact that the two dimensions remain analytically
distinguishable.

\subsubsection{A Provisional Co-Applicability Schema}

For the purposes of a first formal reconstruction,
consider the following schema:

\[
\mathsf{OI}:
\quad
\forall x \in D_E
\left[
E(x)
\rightarrow
\left(B(x)\land I(x)\right)
\right]
\]

The schema states that, for every existent in the
relevant domain, both dimensions apply.

It is intended to exclude a representation in which
an existent is ontologically established as Being
while Identity is wholly absent, or in which Identity
is treated as an independently constituted actuality
without the Being of the existent.

However, the schema requires an important qualification.
The conjunction \(B(x)\land I(x)\) represents
co-applicability. It does not, by itself, capture
the full meaning of ontological inseparability.

In particular, co-applicability alone does not
establish the ontological relation between the
dimensions. It must therefore be supplemented
by an explicit interpretation of inseparability.

Let \(\mathsf{Sep}_O(B,I)\) denote the possibility
of ontological separation between the two dimensions
within the relevant framework. The commitment may
then be represented at the meta-level as:

\[
\mathsf{OI}:
\quad
\neg \mathsf{Sep}_O(B,I)
\]

This expression is not intended to introduce
a fully specified theory of ontological separation.
It records the constraint that the framework
does not admit an interpretation in which
Being and Identity are independently constituted
as separable ontological conditions.

The two expressions therefore perform different
tasks:

\[
\forall x \in D_E
\left[
E(x)\rightarrow(B(x)\land I(x))
\right]
\]

expresses joint applicability, whereas

\[
\neg \mathsf{Sep}_O(B,I)
\]

expresses the ontological constraint against
their separation.

The precise semantics of \(\mathsf{Sep}_O\)
must be specified before the second expression
can function as a fully interpreted formula
in a formal system.

\subsubsection{Inseparability Without Conceptual Identity}

Ontological Inseparability must not be formalized
as the identity of Being and Identity.

In particular, the following expression is
not an adequate representation of the commitment:

\[
B(x) \leftrightarrow I(x)
\]

A biconditional of this kind may express
coextension under a particular interpretation.
It does not, by itself, express ontological
inseparability while preserving analytical
distinction.

Depending on the semantics assigned to the
predicates, it may instead suggest that
Being and Identity have exactly the same
extension or perform the same conceptual role.

That result would exceed the intended commitment.

The present framework requires a representation
in which both dimensions apply to every
existent while remaining distinguishable
in their analytical contents.

The distinction can therefore be expressed
as a conjunction of independent requirements:

\[
\mathsf{AD}
\land
\mathsf{OI}
\]

The conjunction records that analytical
distinction and ontological inseparability
are jointly adopted. It does not reduce
one to the other.

\subsection{Formalization of Absence of Ontological Priority}

\subsubsection{The Content of the Commitment}

Absence of Ontological Priority states that
neither Being nor Identity has ontological
precedence over the other.

The commitment addresses a question different
from analytical distinction and inseparability.

A framework might recognize that two dimensions
are analytically distinguishable and ontologically
inseparable while nevertheless claiming that
one has explanatory or ontological priority.

The present framework rejects such an ordering
between Being and Identity.

The formalization must therefore distinguish
the absence of priority from the absence of
any relation between the dimensions.

\subsubsection{A Provisional Priority Relation}

Let \(\prec_O\) denote a proposed relation
of ontological priority between the relevant
dimensions.

The expression

\[
B \prec_O I
\]

indicates that Being has ontological priority
over Identity.

Likewise,

\[
I \prec_O B
\]

indicates that Identity has ontological
priority over Being.

Absence of Ontological Priority may then
be represented as:

\[
\mathsf{AP}:
\quad
\neg(B\prec_O I)
\land
\neg(I\prec_O B)
\]

The schema excludes both directions of
ontological priority.

Its interpretation must remain limited
to ontological priority. It does not
exclude every possible relation between
Being and Identity, nor does it establish
that they are equivalent or unrelated.

\subsubsection{Priority and Other Relations}

The relation \(\prec_O\) must not be confused
with every form of conceptual, explanatory,
logical, temporal, causal, or epistemic
ordering.

A concept may be introduced before another
in an exposition without being ontologically
prior to it.

A proposition may be used as a premise
in an argument without the corresponding
ontological dimension being prior in reality.

A phenomenon may be recognized or described
before another phenomenon without this
order determining the ontological structure
of the things involved.

The present commitment concerns ontological
priority specifically.

Consequently, the following inference is
invalid:

\[
\neg(B\prec_O I)
\land
\neg(I\prec_O B)
\quad \therefore \quad
B = I
\]

The rejection of priority does not establish
identity or equivalence.

Nor does it establish that Being and Identity
have no relations other than inseparability.
The framework leaves room for further
analysis of their relation, provided that
such analysis does not reintroduce an
unsupported ontological hierarchy.

\subsection{The Joint Formal Architecture}

The three commitments must be represented
as distinct components of the foundational
architecture.

Let:

\[
\begin{aligned}
\mathsf{AD} &:
\text{Analytical Distinction},\\
\mathsf{OI} &:
\text{Ontological Inseparability},\\
\mathsf{AP} &:
\text{Absence of Ontological Priority}.
\end{aligned}
\]

Their joint adoption may be represented
schematically as:

\[
\mathsf{F}_{BI}
=
\mathsf{AD}
\land
\mathsf{OI}
\land
\mathsf{AP}
\]

The expression \(\mathsf{F}_{BI}\) denotes
the combined foundational position concerning
Being and Identity.

It must not be interpreted as a proof
that one commitment follows from the other
two. It records their joint adoption
within the present framework.

The logical architecture can be summarized
as follows:

\[
\begin{array}{c}
\mathsf{AD}\\[4pt]
\mathsf{OI}\\[4pt]
\mathsf{AP}
\end{array}
\quad
\Longrightarrow_{\text{joint adoption}}
\quad
\mathsf{F}_{BI}
\]

The arrow is explanatory rather than
a deductive inference. It indicates
that the three commitments together
constitute the adopted foundational
position.

A formal system may subsequently adopt
their representations as separate axioms.
If it does so, each axiom must retain
its own content and interpretation.

\subsection{The Difference Between Joint Adoption and Logical Derivation}

The conjunction of the three commitments
must not be mistaken for a derivation
of any one commitment from the others.

The formal reconstruction must therefore
distinguish two relations:

\begin{itemize}
    \item \textbf{Joint adoption:} multiple
    commitments are accepted together
    as part of the foundational position.

    \item \textbf{Logical derivation:} a claim
    follows from specified premises under
    specified rules of inference.
\end{itemize}

The present section establishes the
first relation. It does not establish
the second.

To claim that Analytical Distinction
follows from Ontological Inseparability
and Absence of Ontological Priority,
for example, would require a derivation
within a specified formal language.

The same applies to every other proposed
dependency among the three commitments.

Their mutual relations must be investigated
through formal consistency and independence
analysis rather than assumed from their
philosophical association.

\subsection{Formal Non-Equivalence and the Preservation of Distinctions}

The three foundational commitments
must be protected against several
unwarranted substitutions.

First, analytical distinction is not
ontological separation.

Second, ontological inseparability
is not conceptual identity.

Third, absence of ontological priority
is not the absence of every possible
relation.

Fourth, joint adoption is not logical
derivation.

These distinctions can be represented
as constraints on interpretation:

\[
\mathsf{AD}
\not\equiv
\mathsf{OI}
\]

\[
\mathsf{OI}
\not\equiv
(B=I)
\]

\[
\mathsf{AP}
\not\equiv
\neg\mathsf{Relation}(B,I)
\]

\[
\mathsf{F}_{BI}
\not\equiv
\text{a derivation of each of its components}
\]

The symbol \(\not\equiv\) is used here
to indicate non-equivalence of theoretical
commitments, not to introduce a complete
formal semantics of equivalence.

These constraints are important because
the choice of formal symbols can make
different philosophical claims appear
similar when their contents are not
the same.

The purpose of formal reconstruction
is to expose such differences rather
than conceal them.

\subsection{The Risk of Reifying Formal Expressions}

Formal languages often represent
concepts through predicates, constants,
functions, relations, and variables.

These devices are useful for expressing
logical structures, but their use
can introduce unintended ontological
assumptions.

If Being and Identity are represented
as two constants in an object language,
the representation may suggest that
they are two objects belonging to
the same domain.

If they are represented as properties,
the formalization may suggest that
Identity is a property added to
an independently constituted existent.

If they are represented as two
independent domains, the representation
may suggest that each can be specified
without the other.

None of these interpretations is
licensed merely by the foundational
commitments of Part I.

The present formalization therefore
uses meta-level expressions to
represent the distinction and
relations between Being and Identity.

A more developed formal language
may introduce additional structures,
but it must explain why those
structures preserve the intended
ontological interpretation.

The choice of formal representation
is consequently not philosophically
neutral. It must be assessed according
to the commitments it makes explicit
and the assumptions it may inadvertently
introduce.

\subsection{The Limits of the Present Formalization}

The schemas developed in this section
are sufficient to identify the logical
requirements of the three commitments.
They are not yet sufficient to establish
a complete formal theory of Being and
Identity.

Several matters remain unresolved.

\subsubsection{The Semantics of Ontological Applicability}

The formal expressions \(B(x)\) and
\(I(x)\) require a more precise semantic
interpretation.

The framework must determine how
these indicators can be used without
treating Being and Identity as
independent properties of an existent.

The issue cannot be resolved merely
by choosing a different symbol.
It requires an account of the
ontological role of the dimensions
being represented.

\subsubsection{The Semantics of Inseparability}

The expression \(\mathsf{Sep}_O(B,I)\)
requires a formal account of what
counts as ontological separation.

Such an account must distinguish
ontological separation from analytical
distinction, conceptual abstraction,
formal decomposition, and explanatory
differentiation.

Without this distinction, the
formalization may either become
too weak to express the commitment
or too strong by excluding forms
of analysis that the framework
explicitly permits.

\subsubsection{The Meaning of Ontological Priority}

The relation \(\prec_O\) requires
further specification.

Ontological priority may involve
different proposed forms of
dependence, grounding, constitution,
or explanatory precedence.

These relations must not be treated
as interchangeable without argument.

The present schema rejects priority
between Being and Identity at the
level of the foundational commitment.
It does not yet provide a complete
theory of every possible relation
that might be classified as priority.

\subsubsection{The Formal Independence of the Commitments}

The three commitments must be
examined to determine whether
any one can be derived from the
others under a specified formal
interpretation.

This question requires explicit
models, alternative interpretations,
or derivation analysis.

The fact that the commitments
are conceptually distinct does
not by itself prove their formal
independence.

\subsubsection{The Appropriate Formal Language}

The present reconstruction has
not committed to a final single-sorted,
many-sorted, predicate-based,
relational, or other formal language.

The suitability of these alternatives
depends upon the requirements of
the ontology and the interpretive
risks associated with each.

The selection must be made only
after the relevant formal domains,
types, and semantic conditions
have been sufficiently specified.

\subsection{Concluding Principle}

The formal reconstruction of the
principles of identity requires
three distinct representations.

Analytical Distinction concerns
the difference between the
conceptual roles through which
Being and Identity are examined.

Ontological Inseparability concerns
their inseparable status in the
actuality of an existent.

Absence of Ontological Priority
concerns the rejection of an
ontological ordering in which
either dimension is prior to
the other.

Their joint adoption may be
represented as a foundational
architecture, but their
individual contents must remain
independent and available for
separate examination.

The governing principle of
this formalization is:

\begin{quote}
A formal representation of
Being and Identity must preserve
their analytical distinction,
their ontological inseparability,
and the absence of ontological
priority, without converting
any one of these commitments
into an assertion of conceptual
identity, ontological separation,
or an unsupported relation
of dependence.
\end{quote}

The next section examines the
formal propositions and derived
claims that may be developed
from these commitments. It
must distinguish what is
adopted as a starting point
from what can be established
through explicit inference.

% =============================================================================
% END OF SECTION 06 OF CHAPTER 07
% =============================================================================